% ============================================================================
% Chapter 2, Section 1: Introduction
% ============================================================================

\section{Background}\label{sec:background}

\subsection{Introduction}\label{sec:introduction}

Partial observability, long-horizon planning, and multimodal perception are core challenges in sequential decision-making, yet few existing benchmarks stress all three simultaneously under realistic conditions. Standard testbeds tend to isolate one axis: imperfect-information games emphasize equilibrium computation in short episodes, while open-ended environments test exploration but lack adversarial opponents or extended temporal commitments. As frontier language and vision--language models (LLMs/VLMs) are increasingly deployed as autonomous agents---navigating codebases, operating robotic systems, and interacting with complex virtual environments---the gap between what existing benchmarks measure and what real-world autonomy demands has grown acute. The field needs evaluation frameworks that natively test the kinds of sustained, multimodal reasoning that these agents must exhibit.

Pok\'{e}mon role-playing games (RPGs) offer a compelling environment for this purpose. They combine a vast, partially observable state space with long-horizon narrative progression that demands thousands of cumulative decisions spanning exploration, resource management, combat, and spatial navigation. The recent surge of interest in Pok\'{e}mon-playing AI systems---Claude Plays Pok\'{e}mon~\citep{anthropic2025visible}, Gemini Plays Pok\'{e}mon~\citep{zhang2025geminiplayspokemon}, and GPT-5's completion of Pok\'{e}mon Red~\citep{engadget2025gptplayspokemon}---has demonstrated both the appeal of this domain and the limitations of current approaches. These efforts were fragmented across different games, harnesses, and evaluation criteria, making it impossible to disentangle whether success should be attributed to the underlying model, the agent architecture, or hardcoded assumptions that simplified perception.

This thesis addresses the fragmentation problem and advances the state of the art across three complementary dimensions:

\begin{enumerate}[leftmargin=*]
    \item \textbf{A live speedrunning benchmark} for testing frontier model capabilities in long-duration, embodied settings. Originally developed as the Speedrunning Track of the \pokeagent{} Challenge at NeurIPS 2025~\citep{karten2025pokeagent}, the benchmark formalizes Pok\'{e}mon Emerald gameplay as a structured evaluation with standardized milestones, quantitative metrics, and infrastructure for fair cross-paradigm comparison. We have since extended it from the tutorial segment (through the first gym) to a long-horizon evaluation spanning three gym leaders, where experiments routinely exceed twelve hours of continuous gameplay.

    \item \textbf{A state-of-the-art custom VLM-agent harness} with abstractions for coherence, planning, context management, and subagent orchestration. The PokeAgent baseline evolved from a simple four-module scaffold during the NeurIPS competition into a multi-agent system with a trajectory window for session-level context, hierarchical memory (short-term and long-term), a skill library, delegated subagents, and sandboxed code execution. Comparisons against commercial coding agents (Claude Code, Gemini CLI, Codex CLI) demonstrate that the harness---not the model---is the primary bottleneck for embodied agent performance.

    \item \textbf{AutoEvolve}, a meta-harness framework that automates experiential learning to evolve agent scaffolding mid-episode without environment resets. Starting from a minimal interface (frame observations and button inputs), AutoEvolve bootstraps its own system prompt, creates specialized subagents, builds a skill library, and refines all of these through in-context optimization over trajectory segments. Unlike prior prompt-evolution methods that require full episodes and environment resets~\citep{agrawal2025gepa}, AutoEvolve operates in a fully online, reset-free manner---enabling continuous self-improvement during a single extended gameplay session.
\end{enumerate}

The remainder of this thesis is organized as follows. The present chapter provides motivation and context for the work (Section~\ref{sec:motivation}), surveys related research across embodied agents, self-adaptation, and the emerging field of meta-harnesses (Section~\ref{sec:related_work}), and establishes the formal framework used throughout subsequent chapters (Section~\ref{sec:formalization}). Chapter~3 describes the speedrunning benchmark in detail: its environment design, evaluation criteria, and baselines. Chapter~4 traces the development of the PokeAgent harness from its competition-era origins through its current multi-agent architecture, presenting early performance results that motivate the need for automated harness evolution. Chapter~5 presents AutoEvolve---the primary research contribution of this thesis---including its methodology, experimental design across three scaffold tiers, and results. Chapter~6 concludes with a discussion of implications and future directions, including model distillation, full-game evaluations, and extensions to domains beyond Pok\'{e}mon.

We note that Chapters~3 and~4 draw substantially from the published \pokeagent{} Challenge benchmark paper~\citep{karten2025pokeagent} and from our midterm report, while Chapter~5 builds on the ongoing AutoEvolve collaboration. Throughout, we focus on the speedrunning domain within Pok\'{e}mon Emerald; competitive battling, while an important component of the broader \pokeagent{} Challenge, falls outside the scope of this thesis except where it provides necessary context.
